\documentclass[a4paper, 12pt]{article} % Formato de la plantilla


\usepackage{titlesec}
\usepackage[utf8]{inputenc} % UTF8 Encoding
\usepackage[margin=2cm, top=2cm, includefoot]{geometry} % Geometría
\usepackage{graphicx} % Inserción de imagenes
\usepackage[table,xcdraw]{xcolor} % Detección de colores
\usepackage[most]{tcolorbox} % Inserción de cuadros en la portada
\usepackage{fancyhdr} % Estilo de la página
\usepackage[hidelinks]{hyperref} % Gestión de hipervínculos
\usepackage{listings} % Inserción de código
\usepackage{parskip} % Tabulación del documento
\usepackage[figurename=Figure]{caption} % Para cambiar el nombre del caption de las fotos
\usepackage{smartdiagram} % Inserción de diagramas
\usepackage{zed-csp} % Inserción de esquemas
\usepackage{float} % Para la colocación de imagenes y tablas, sirve para H
\usepackage{colortbl} % Colores de tablas
\usepackage{tabularx} % Mejores tablas
%\usepackage{bidi} % Bidirectional typsetting bugs of XeTeX
\usepackage{apacite}
\usepackage{anyfontsize}

% Define Colors

% Colors from color picker: https://www.google.com/search?q=color+picker

\definecolor{green}{HTML}{69A84F}
\definecolor{gray}{HTML}{B8B6B4}
\definecolor{yellow}{HTML}{acba0f}
\definecolor{red}{HTML}{f70000}
\definecolor{blue}{HTML}{0370ff}
\definecolor{pink}{HTML}{ff03d1}

% Font

\renewcommand{\familydefault}{\sfdefault}

% Tables 

% Table generator: https://www.latex-tables.com/

\renewcommand{\contentsname}{Table of Contents}
\renewcommand{\arraystretch}{1.5}
\renewcommand{\tabularxcolumn}[1]{>{\small}m{#1}}

% Header & Footer

\setlength{\headheight}{40.2pt}
\pagestyle{fancy}
\fancyhf{}
\lhead{\includegraphics[height=0.5cm]{media/avatar/avatar.png}}
\rhead{Application Security Review}
\renewcommand{\headrulewidth}{1pt}
\renewcommand{\headrule}{\hbox to\headwidth{\color{gray}\leaders\hrule height \headrulewidth\hfill}}


\lfoot{CompanyName Confidential}
\cfoot{Copyright © CompanyName}
\rfoot{Page \thepage}
\renewcommand{\footrulewidth}{1pt}
\renewcommand{\footrule}{\hbox to\headwidth{\color{gray}\leaders\hrule height 
\footrulewidth\hfill}}

% Hypertext

\hypersetup{
%    colorlinks=true, Enable Colors for Links
%    linkcolor=blue, All Links will have the color declared
    filecolor=magenta,      
    urlcolor=cyan
}

% Sections

\setcounter{secnumdepth}{5}

\titleformat{\paragraph}
{\normalfont\normalsize\bfseries}{\theparagraph}{1em}{}
\titlespacing*{\paragraph}
{0pt}{3.25ex plus 1ex minus .2ex}{1.5ex plus .2ex}

% Colors for Code Listings

\definecolor{codegreen}{rgb}{0,0.6,0}
\definecolor{codegray}{rgb}{0.5,0.5,0.5}
\definecolor{codepurple}{rgb}{0.58,0,0.82}
\definecolor{backcolour}{rgb}{0.95,0.95,0.92}

\lstdefinestyle{mystyle}{
    backgroundcolor=\color{backcolour},   
    commentstyle=\color{codegreen},
    keywordstyle=\color{magenta},
    numberstyle=\tiny\color{codegray},
    stringstyle=\color{codepurple},
    basicstyle=\ttfamily\footnotesize,
    breakatwhitespace=false,         
    breaklines=true,                 
    captionpos=b,                    
    keepspaces=true,                 
    numbers=left,                    
    numbersep=5pt,                  
    showspaces=false,                
    showstringspaces=false,
    showtabs=false,                  
    tabsize=2
}

\lstset{style=mystyle}

% Start of the Document

\begin{document}
    % Cover Page
    \begin{titlepage}
        \centering
        \vfill
        \begin{figure}[t]
            \centering
            \includegraphics[width=\textwidth]{media/avatar/avatar.png}
        \end{figure}
        \vfill
        \begin{center}
            \vspace*{\fill}
            {\fontsize{20}{10}\selectfont\textbf{Client Company Name}}\par
            {\fontsize{20}{10}\selectfont\textbf{Security Assessment Report}}\par
            \vspace{5cm}
            \Huge Business Confidential\par
            \vspace{1cm}
            \begin{flushleft}
            \small\textit{Date: May, 17, 2021}\par
            \small\textit{Project: 898-21}\par
            \small\textit{Version: 1.0}\par
            \end{flushleft}
            \vspace{5pt}
            \vfill
        \end{center}
    \end{titlepage}

    \clearpage 

    % ---------------------
    % Table of Contents
    \clearpage
    \tableofcontents
    \clearpage
    % ----------------------


    \section{Confidentiality Statement}

    This document is the exclusive property of Client Company and CompanyName). This document contains proprietary and confidential information. Duplication, redistribution, or use, in whole or in part, in any form, requires consent of both Client Company and CompanyName.

    CompanyName may share this document with auditors under non-disclosure agreements to demonstrate penetration test requirement compliance.

    \section{Disclaimer}

    A penetration test is considered a snapshot in time. The findings and recommendations reflect the information gathered during the assessment and not any changes or modifications made outside of that period. Time-limited engagements do not allow for a full evaluation of all security controls. CompanyName prioritized the assessment to identify the weakest security controls an attacker would exploit. CompanyName recommends conducting similar assessments on an annual basis by internal or third-party assessors to ensure the continued success of the controls.

    \section{Contact Information}


    \begin{table}[H]
        \centering
        \begin{tabular}{|>{\raggedright\arraybackslash}m{23mm}|m{60mm}|m{60mm}|} 
        \hline
        \rowcolor[rgb]{0.498,0.498,0.498} Name & Title                           & Contact Information                                       \\ 
        \hline
        \multicolumn{3}{|l|}{{\cellcolor[rgb]{0.851,0.851,0.851}}Client Company}                                                             \\ 
        \hline
        John Smith                             & VP, Information Security (CISO) & Office: (555) 555-5555\newline
        Email: john.smith@demo.com       \\ 
        \hline
        Jim Smith                              & IT Manager                      & Office: (555) 555-5555\newline
        Email: jim.smith@demo.com        \\ 
        \hline
        Joe Smith                              & Network Engineer                & Office: (555) 555-5555\newline
        Email: joe.smith@demo.com        \\ 
        \hline
        \multicolumn{3}{|l|}{{\cellcolor[rgb]{0.851,0.851,0.851}}Company Name}                                                               \\ 
        \hline
        John Smith                             & Lead Penetration Tester         & Office: (555) 555-5555\newline
        Email: jsmith@company\_name.com  \\ 
        \hline
        Bob Smith                              & Penetration Tester              & Office: (555) 555-5555\newline
        Email: bsmith@company\_name.com  \\ 
        \hline
        Rob Smith                              & Account Manager                 & Office: (555) 555-5555\newline
        Email: rsmith@company\_name.com  \\
        \hline
        \end{tabular}
        \caption{Contact}
    \end{table}

    \clearpage

    \section{Assessment Overview}

    From June 20th, 2021 to June 29th, 2021, Client Company engaged CompanyName to evaluate the security posture of its infrastructure compared to current industry best practices that included an external penetration test. All testing performed is based on the NIST SP 800-115 Technical Guide to Information Security Testing and Assessment, OWASP Testing Guide (v4), and customized testing frameworks. 
    Phases of penetration testing activities include the following:
    \begin{itemize}
        \item   Planning – Customer goals are gathered and rules of engagement obtained.
        \item	Discovery – Perform scanning and enumeration to identify potential vulnerabilities, weak areas, and exploits.
        \item	Attack – Confirm potential vulnerabilities through exploitation and perform additional discovery upon new access.
        \item	Reporting – Document all found vulnerabilities and exploits, failed attempts, and company strengths and weaknesses.
    \end{itemize}
    
    \subsection{Assessment Components}

    \textbf{External Penetration Test}

    An external penetration test emulates the role of an attacker attempting to gain access to an internal network without internal resources or inside knowledge. A CompanyName engineer attempts to gather sensitive information through open-source intelligence (OSINT), including employee information, historical breached passwords, and more that can be leveraged against external systems to gain internal network access. The engineer also performs scanning and enumeration to identify potential vulnerabilities in hopes of exploitation.

    \clearpage

    \section{Finding Severity Ratings}

    The following table defines levels of severity and corresponding CVSS score range that are used throughout the document to assess vulnerability and risk impact.

    % \usepackage{colortbl}


    \begin{table}[H]
        \centering
        \begin{tabular}{|>{\raggedright\arraybackslash}m{23mm}|m{40mm}|m{90mm}|}
        \hline
        \rowcolor[rgb]{0.502,0.502,0.502} Severity    & CVSS V3 Score Range & Definition                                                                                                                                                                                               \\ 
        \hline
        {\cellcolor[rgb]{0.753,0,0}}Critical          & 9.0-10.0            & Exploitation is straightforward and usually results in system-level
        compromise.~ It is advised to form a
        plan of action and patch immediately.                                                       \\ 
        \hline
        {\cellcolor{red}}High                         & 7.0-8.9             & Exploitation is more difficult but could cause elevated privileges
        and potentially a loss of data or downtime.~
        It is advised to form a plan of action and patch as soon as possible.                \\ 
        \hline
        {\cellcolor[rgb]{1,0.753,0}}Moderate          & 4.0-6.9             & Vulnerabilities exist but are not exploitable or require extra steps
        such as social engineering.~ It is
        advised to form a plan of action and patch after high-priority issues have
        been resolved.  \\ 
        \hline
        {\cellcolor{yellow}}Low                       & 0.1-3.9             & Vulnerabilities are non-exploitable but would reduce an
        organization’s attack surface.~ It is
        advised to form a plan of action and patch during the next maintenance
        window.                       \\ 
        \hline
        {\cellcolor[rgb]{0,0.439,0.753}}Informational & N/A                 & No vulnerability exists.~
        Additional information is provided regarding items noticed during
        testing, strong controls, and additional documentation.                                                  \\
        \hline
        \end{tabular}
        \caption{Severity Ratings}
    \end{table}

    \clearpage

    \section{Scope}
    
    \begin{table}[H]
        \centering
        \begin{tabular}{!{\color{black}\vrule}l|l!{\color{black}\vrule}} 
        \hline
        \rowcolor[rgb]{0.498,0.498,0.498} Assessment & Details                            \\ 
        \hline
        External Penetration Test                    & 192.168.0.0/24, 
        192.168.1.0/24  \\
        \hline
        \end{tabular}
        \caption{Scope}
    \end{table}

    \subsection{Scope Exclusions}

    Per client request, CompanyName did not perform any X attacks during testing.

    Client Allowances

    ClientCompany did not provide any allowances to assist the testing.

    \section{Executive Summary}

    CompanyName evaluated ClientCompany’s external security posture through an external network penetration test from June 20th, 2019 to June 29th, 2019.  By leveraging a series of attacks, CompanyName found critical level vulnerabilities that allowed full internal network access to the ClientCompany headquarter office. It is highly recommended that Client Company address these vulnerabilities as soon as possible as the vulnerabilities are easily found through basic reconnaissance and exploitable without much effort.

    \section{Attack Summary}

    The following table describes how CompanyName gained internal network access, step by step:

    Lorem ipsum

    \subsection{Security Strengths}

    SIEM alerts of vulnerability scans

    During the assessment, lorem ipsum dolor sit amet, consectetur adipiscing elit. Sed scelerisque sapien id risus tristique, vel molestie erat dapibus. Maecenas eleifend nunc sit amet purus aliquam, eu placerat sem dapibus. Orci varius natoque penatibus et magnis dis parturient montes, nascetur ridiculus mus. Pellentesque rhoncus commodo ultricies. Mauris quis quam ut tellus finibus lobortis. Donec dapibus leo a consectetur mollis. Donec nec mi euismod, vehicula tortor eu, feugiat massa. Curabitur at risus sed lorem fringilla auctor non pellentesque sem. Suspendisse urna arcu, sodales in risus sed, interdum accumsan mauris. 

    \subsection{Security Weaknesses}

    \subsubsection{Missing Multi-Factor Authentication}

    Lorem ipsum dolor sit amet, consectetur adipiscing elit. Sed scelerisque sapien id risus tristique, vel molestie erat dapibus. Maecenas eleifend nunc sit amet purus aliquam, eu placerat sem dapibus. Orci varius natoque penatibus et magnis dis parturient montes, nascetur ridiculus mus. Pellentesque rhoncus commodo ultricies. Mauris quis quam ut tellus finibus lobortis. Donec dapibus leo a consectetur mollis. Donec nec mi euismod, vehicula tortor eu, feugiat massa. Curabitur at risus sed lorem fringilla auctor non pellentesque sem. Suspendisse urna arcu, sodales in risus sed, interdum accumsan mauris. 

    \subsubsection{Weak Password Policy}

    Lorem ipsum dolor sit amet, consectetur adipiscing elit. Sed scelerisque sapien id risus tristique, vel molestie erat dapibus. Maecenas eleifend nunc sit amet purus aliquam, eu placerat sem dapibus. Orci varius natoque penatibus et magnis dis parturient montes, nascetur ridiculus mus. Pellentesque rhoncus commodo ultricies. Mauris quis quam ut tellus finibus lobortis. Donec dapibus leo a consectetur mollis. Donec nec mi euismod, vehicula tortor eu, feugiat massa. Curabitur at risus sed lorem fringilla auctor non pellentesque sem. Suspendisse urna arcu, sodales in risus sed, interdum accumsan mauris. 

    \subsubsection{Unrestricted Logon Attempts}

    Lorem ipsum dolor sit amet, consectetur adipiscing elit. Sed scelerisque sapien id risus tristique, vel molestie erat dapibus. Maecenas eleifend nunc sit amet purus aliquam, eu placerat sem dapibus. Orci varius natoque penatibus et magnis dis parturient montes, nascetur ridiculus mus. Pellentesque rhoncus commodo ultricies. Mauris quis quam ut tellus finibus lobortis. Donec dapibus leo a consectetur mollis. Donec nec mi euismod, vehicula tortor eu, feugiat massa. Curabitur at risus sed lorem fringilla auctor non pellentesque sem. Suspendisse urna arcu, sodales in risus sed, interdum accumsan mauris. 

    \subsubsection{Credentials Reuse} 

    Lorem ipsum dolor sit amet, consectetur adipiscing elit. Sed scelerisque sapien id risus tristique, vel molestie erat dapibus. Maecenas eleifend nunc sit amet purus aliquam, eu placerat sem dapibus. Orci varius natoque penatibus et magnis dis parturient montes, nascetur ridiculus mus. Pellentesque rhoncus commodo ultricies. Mauris quis quam ut tellus finibus lobortis. Donec dapibus leo a consectetur mollis. Donec nec mi euismod, vehicula tortor eu, feugiat massa. Curabitur at risus sed lorem fringilla auctor non pellentesque sem. Suspendisse urna arcu, sodales in risus sed, interdum accumsan mauris. 

    \clearpage

    \section{External Penetration Test Findings}

    \textbf{Insufficient Lockout Policy - Outlook Web App}
    \begin{table}[H]
        \centering
        \begin{tabular}{|>{\raggedright\arraybackslash}m{20mm}|m{130mm}|m{60mm}|} 
        \hline
        {\cellcolor[rgb]{0.816,0.808,0.808}}Description: & DC allowed unlimited logon attempts against their Outlook
          Web App (OWA) services. This configuration allowed brute force and password
          guessing attacks in which CompanyName used to gain access to ClientCompany internal network.  \\ 
        \hline
        {\cellcolor[rgb]{0.816,0.808,0.808}}Impact:      & Critical                                                                                                                                                                                                                \\ 
        \hline
        {\cellcolor[rgb]{0.816,0.808,0.808}}System:      & 192.168.0.5                                                                                                                                                                                                             \\ 
        \hline
        {\cellcolor[rgb]{0.816,0.808,0.808}}References:  & NIST SP800-53r4 AC-17 - Remote Access\newline
          NIST
          SP800-53r4 AC-7(1) - Unsuccessful Logon Attempts                                                                                     \\
        \hline
        \end{tabular}
        \caption{Vulnerability name}
    \end{table}

    \textbf{Exploitation Proof of Concept}

    Lorem ipsum dolor sit amet, consectetur adipiscing elit. Sed scelerisque sapien id risus tristique, vel molestie erat dapibus. Maecenas eleifend nunc sit amet purus aliquam, eu placerat sem dapibus. Orci varius natoque penatibus et magnis dis parturient montes, nascetur ridiculus mus. Pellentesque rhoncus commodo ultricies. Mauris quis quam ut tellus finibus lobortis. Donec dapibus leo a consectetur mollis. Donec nec mi euismod, vehicula tortor eu, feugiat massa. Curabitur at risus sed lorem fringilla auctor non pellentesque sem. Suspendisse urna arcu, sodales in risus sed, interdum accumsan mauris. 

    \begin{figure}[H]
        \centering
        \includegraphics[width=\textwidth]{media/img/test.png}
        \caption{Vulnerability Name Exploit in Action Phase 1}
    \end{figure}


    Lorem ipsum dolor sit amet, consectetur adipiscing elit. Sed scelerisque sapien id risus tristique, vel molestie erat dapibus. Maecenas eleifend nunc sit amet purus aliquam, eu placerat sem dapibus. Orci varius natoque penatibus et magnis dis parturient montes, nascetur ridiculus mus. Pellentesque rhoncus commodo ultricies. Mauris quis quam ut tellus finibus lobortis. Donec dapibus leo a consectetur mollis. Donec nec mi euismod, vehicula tortor eu, feugiat massa. Curabitur at risus sed lorem fringilla auctor non pellentesque sem. Suspendisse urna arcu, sodales in risus sed, interdum accumsan mauris. 

    \begin{figure}[H]
        \centering
        \includegraphics[width=\textwidth]{media/img/test.png}
        \caption{Vulnerability Name Exploit in Action Phase 2}
    \end{figure}


    Lorem ipsum dolor sit amet, consectetur adipiscing elit. Sed scelerisque sapien id risus tristique, vel molestie erat dapibus. Maecenas eleifend nunc sit amet purus aliquam, eu placerat sem dapibus. Orci varius natoque penatibus et magnis dis parturient montes, nascetur ridiculus mus. Pellentesque rhoncus commodo ultricies. Mauris quis quam ut tellus finibus lobortis. Donec dapibus leo a consectetur mollis. Donec nec mi euismod, vehicula tortor eu, feugiat massa. Curabitur at risus sed lorem fringilla auctor non pellentesque sem. Suspendisse urna arcu, sodales in risus sed, interdum accumsan mauris. 

    \begin{figure}[H]
        \centering
        \includegraphics[width=\textwidth]{media/img/test.png}
        \caption{Vulnerability Name Exploit in Action Result}
    \end{figure}

    Lorem ipsum dolor sit amet, consectetur adipiscing elit. Sed scelerisque sapien id risus tristique, vel molestie erat dapibus. Maecenas eleifend nunc sit amet purus aliquam, eu placerat sem dapibus. Orci varius natoque penatibus et magnis dis parturient montes, nascetur ridiculus mus. Pellentesque rhoncus commodo ultricies. Mauris quis quam ut tellus finibus lobortis. Donec dapibus leo a consectetur mollis. Donec nec mi euismod, vehicula tortor eu, feugiat massa. Curabitur at risus sed lorem fringilla auctor non pellentesque sem. Suspendisse urna arcu, sodales in risus sed, interdum accumsan mauris. 

    \textbf{Remediation}

    Lorem ipsum dolor sit amet, consectetur adipiscing elit. Sed scelerisque sapien id risus tristique, vel molestie erat dapibus. Maecenas eleifend nunc sit amet purus aliquam, eu placerat sem dapibus. Orci varius natoque penatibus et magnis dis parturient montes, nascetur ridiculus mus. Pellentesque rhoncus commodo ultricies. Mauris quis quam ut tellus finibus lobortis. Donec dapibus leo a consectetur mollis. Donec nec mi euismod, vehicula tortor eu, feugiat massa. Curabitur at risus sed lorem fringilla auctor non pellentesque sem. Suspendisse urna arcu, sodales in risus sed, interdum accumsan mauris. 

\end{document}